\documentclass[../main.tex]{subfiles}

\begin{document}
\chapter{Diseño}
\label{C:disenno}


\section{Diseño del sistema}

\subsection{Arquitectura}

\subsubsection{Modelo}
\subsubsection{Vista}
\subsubsection{Controlador}

\subsection{Diseño de vistas}


\clearpage

\subsubsection{Configuración}

La pantalla de configuración es una de las pantallas más importantes de la aplicación, ya que desde esta pantalla controlamos todos los aspectos personalizables de la aplicación, como el aspecto, idioma o incluso los parámetros para el compilador.

Con el motivo de la multifuncionalidad de esta vista, se ha diseñado con un sistema de pestañas y sub-pestañas, en el que cada pestaña es un aspecto de la aplicación, básicamente lo dividiríamos en tres pestañas principales diferenciadas: idioma, compilador y plugins.

\paragraph{Configuración de idioma}

La configuración de idioma permite al usuario seleccionar el idioma de la interfaz de la aplicación, ofreciendo soporte multiidioma para una mejor experiencia de usuario.

\begin{figure}[!htb]
    \centering
    \includegraphics[width=0.8\textwidth]{Vistas/Dialogo_Config_idioma.png}
    \caption{Diálogo de configuración de idioma.}
    \label{fig:config_idioma}
\end{figure}

\clearpage

\paragraph{Configuración de compiladores}

El sistema permite configurar diferentes toolchains de compilación según el target de desarrollo. Se incluyen configuraciones específicas para GBDK-2020 y Z88DK.

\begin{figure}[!htb]
    \centering
    \includegraphics[width=0.8\textwidth]{Vistas/Dialogo_Config_Compilador_GBDK.png}
    \caption{Diálogo de configuración del compilador GBDK-2020.}
    \label{fig:config_gbdk}
\end{figure}

\begin{figure}[!htb]
    \centering
    \includegraphics[width=0.8\textwidth]{Vistas/Dialogo_Config_Compilador_z88dk.png}
    \caption{Diálogo de configuración del compilador Z88DK.}
    \label{fig:config_z88dk}
\end{figure}

\paragraph{Configuración de plugins}

La gestión de plugins permite instalar, activar, desactivar y configurar extensiones del sistema, proporcionando flexibilidad y extensibilidad a la aplicación.

\begin{figure}[!htb]
    \centering
    \includegraphics[width=0.8\textwidth]{Vistas/Dialogo_Config_plugins.png}
    \caption{Diálogo de configuración de plugins.}
    \label{fig:config_plugins}
\end{figure}

\subsubsection{Diálogos auxiliares}

El sistema incluye diversos diálogos auxiliares que facilitan la interacción del usuario con la aplicación.

\paragraph{Creación de nuevos proyectos}

\begin{figure}[!htb]
    \centering
    \includegraphics[width=0.8\textwidth]{Vistas/Dialogo_Nuevo_proyecto.png}
    \caption{Diálogo para la creación de nuevos proyectos.}
    \label{fig:nuevo_proyecto}
\end{figure}

\paragraph{Funcionalidades de búsqueda}

El sistema proporciona herramientas de búsqueda tanto en archivos individuales como a nivel de proyecto.

\begin{figure}[!htb]
    \centering
    
\includegraphics[width=0.8\textwidth]{Vistas/Dialogo_Buscar.png}
    \caption{Diálogo de búsqueda en archivos.}
    \label{fig:dialogo_buscar}
\end{figure}

\begin{figure}[!htb]
    \centering
    \includegraphics[width=0.8\textwidth]{Vistas/Dialogo_Buscar_proyecto.png}
    \caption{Diálogo de búsqueda a nivel de proyecto.}
    \label{fig:buscar_proyecto}
\end{figure}

\paragraph{Navegación de código}
 
\begin{figure}[!htb]
    \centering
    \includegraphics[width=0.8\textwidth]{Vistas/Dialogo_Ir_Linea.png}
    \caption{Diálogo para ir a una línea específica.}
    \label{fig:ir_linea}
\end{figure}

\subsubsection{Configuración de emuladores}

La aplicación integra emuladores para probar el software desarrollado. La configuración de JSpeccy permite ajustar los parámetros de emulación del ZX Spectrum.

\begin{figure}[!htb]
    \centering
    \includegraphics[width=0.8\textwidth]{ConfiguracionJSpeccy.png}
    \caption{Interface de configuración del emulador JSpeccy.}
    \label{fig:config_jspeccy}
\end{figure}










\end{document}